% Avsnitt 4.10, ur lectures/L04/README.md.

\section{Sammanfattning}
\label{c4:sec:review}

\needspace{6\baselineskip}
\subsection*{Det här ska du kunna nu}
\begin{itemize}
  \item Förklara vad metastabilitet är, varför en asynkron ingång kan orsaka det, och varför två
    vippor gör det osannolikt snarare än omöjligt.
  \item Tillämpa dubbelvippsynkroniseraren på vilken asynkron ingång som helst, inklusive mönstret
    ”aktivera asynkront, släpp synkront” som används för reset.
  \item Säga varför den kedjan dessutom fungerar som en praktisk men ofullständig studsfiltrering,
    och vad en riktig sådan skulle tillföra.
  \item Ge en modul en \code{generic}, åsidosätta den med en \code{generic map}, skriva en ports
    bredd i termer av den, och bedöma när en generic gör rätt för sig i stället för att läggas till
    av princip.
\end{itemize}

\needspace{6\baselineskip}
\subsection*{Frågor att testa dig själv med}
\begin{enumerate}
  \item Vad händer exakt när en asynkron ingång till en vippa ändras för nära en aktiv klockflank?
  \item Hur hanterar dubbelvippsynkroniseraren det, och varför är det den \emph{andra} vippans
    utgång som är säker att använda?
  \item Varför beter sig kedjan som används här som en studsfiltrering i CircuitVerse men
    studsfiltrerar inte en riktig knapp fullt ut vid \code{50 MHz}, och varför är det en egenskap
    hos klockperioden snarare än hos kretsen? Vad skulle du lägga till för att åtgärda det?
  \item \code{button_sync} har en generic och \code{reset_sync} har ingen. Vad är regeln, och vad
    vore fel med att ge \code{reset_sync} en \code{WIDTH}-generic ”för symmetrins skull”?
\end{enumerate}

\needspace{6\baselineskip}
\subsection*{Vidare läsning}
Erik Pihls videogenomgång av metastabilitet och dubbelvippsynkroniseraren finns på
\url{https://www.youtube.com/watch?v=KrssJRgF13I} som ytterligare, frivillig visning.
\url{https://nandland.com/lesson-13-metastability/} och
\url{https://vhdlwhiz.com/terminology/metastability/} täcker samma begrepp från en något annan
vinkel, användbart om det här kapitlet inte sitter vid första läsningen.

\needspace{6\baselineskip}
\subsection*{Blicken framåt}
\Chapref{c5:ch} går vidare med:
\begin{itemize}
  \item \code{variable}: den enda konstruktion du ännu inte mött, och den skarpa skillnaden mellan
    hur den uppdateras och hur en \code{signal} gör det.
  \item Varför samma ackumulering skriven med en \code{signal} och med en \code{variable} ger två
    olika svar, och vilket av de två som är felet.
  \item Vad verktygskedjan faktiskt bygger: LUT:ar och vippor, vad som begränsar en klockas
    hastighet, och låset du råkar beskriva av misstag.
\end{itemize}
